% Part: many-valued-logic
% Chapter: syntax-and-semantics
% Section: formulas

\documentclass[../../../include/open-logic-section]{subfiles}

\begin{document}

\olfileid[pt]{mvl}{syn}{fml}

\olsection{\usetoken{P}{formula}}

\begin{defn}[Fórmula]
\ollabel{defn:formulas}
O conjunto~$\Frm[L]$ de \emph{!!{formula}s} de uma linguagem
proposicional~$\Lang L$ é definido indutivamente da seguinte maneira:
\begin{enumerate}
\item Toda !!{propositional variable}~$\Obj p_i$ é uma
  !!{formula} atômica.
\item Todo conectivo $0$-ário (constante proposicional) de~$\Lang L$
é uma !!{formula} atômica.
\item Se $\star$ é um conectivo $n$-ário de~$\Lang L$, e $!A_1$,
\dots, $!A_n$ são !!{formula}s, então $\star(!A_1, \dots, !A_n)$ é
  !!a{formula}.
\tagitem{limitClause}{Nada mais é !!a{formula}.}{}
\end{enumerate}
Se $\star$ é $1$-ário, então $\star(!A_1)$ muitas vezes será escrito
simplesmente como $\star !A_1$. Se $\star$ é $2$-ário, $\star(!A_1,!A_2)$
muitas vezes será escrito como $(!A_1 \star !A_2)$. 
\end{defn}

Como de costume, frequentemente omitiremos sem menção os parênteses mais externos.

\begin{ex}
  Na linguagem padrão~$\Lang{L_0}$, $\Obj p_1 \lif (\Obj p_1
  \land \lnot \Obj p_2)$ é uma fórmula. Na linguagem da lógica do produto,
  ela seria escrita, em vez disso, como $\Obj p_1 \lif (\Obj p_1 \odot
  \lnot \Obj p_2)$. Se acrescentarmos à linguagem o conectivo $1$-ário $\triangle$,
  também teremos fórmulas como $\triangle (\Obj p_1
  \land \Obj p_2) \lif (\triangle \Obj p_1 \land \triangle \Obj p_2)$.
\end{ex}

\end{document}
